\documentclass[10pt,a4paper]{article}
\usepackage[T1]{fontenc}
\usepackage[margin=1.7cm]{geometry}
\usepackage{booktabs,tabularx}
\usepackage[hidelinks]{hyperref}
\setlength{\parindent}{0pt}
\setlength{\parskip}{5pt}
\pagestyle{empty}

\begin{document}
\begin{center}
  {\Large\bfseries Market Candles and Weather Forecasts: Data Availability}\\[3pt]
  Latest evidence reviewed: 7 September 2026
\end{center}

\textbf{Market candles.}
The saved Kalshi daily-high-temperature probe finds traded daily candles for
25 of 47 discovered series. It combines live and historical market listings,
selects the highest-volume contract for each event-day, and requests
1440-minute candles. The earliest successful event-dates are below; these are
event labels, not necessarily the UTC dates on which trading began.

\begin{center}
\small
\begin{tabularx}{\linewidth}{@{}lX@{}}
\toprule
First event-date found & Series (each suffix below has the prefix \texttt{KXHIGH}) \\
\midrule
2021-08-06 & \texttt{NY} (New York) \\
2021-08-19 & \texttt{CHI} (Chicago) \\
2023-05-11 & \texttt{AUS, MIA} (Austin, Miami) \\
2024-11-20 & \texttt{DEN, HOU, PHIL} (Denver, Houston, Philadelphia) \\
2025-01-05 & \texttt{LAX} (Los Angeles) \\
2026-01-14 & \texttt{TDC, TLV, TNOLA, TSEA, TSFO} \\
2026-02-04 & \texttt{TATL, TMIN, TPHX} \\
2026-02-05 & \texttt{TBOS} \\
2026-02-11 & \texttt{TDAL, THOU, TOKC, TSATX} \\
2026-08-19 & \texttt{TSAN} \\
2026-08-26 & \texttt{TEWR, TTTN} \\
2026-08-28 & \texttt{TSDF} \\
\bottomrule
\end{tabularx}
\end{center}

The saved results reflect a limit-only search: it stops at the first successful
day per series. Consequently, its \texttt{newest\_event\_day\_with\_candle} field
repeats the first date and \textbf{does not establish the latest available candle}.
The other 22 series have no successful sample in this probe. Neither complete
daily coverage nor availability for every contract or intraday interval was
verified.

\textbf{Weather forecasts.}
The project uses Open-Meteo's Previous Runs API for hourly 2-metre temperature
forecasts at fixed lead times. The provider documents coverage from
\textbf{January 2024 for most models}, with later starts for some models;
\textbf{GFS temperature extends to March 2021}, and JMA GSM/MSM to 2018.
These are provider-stated archive bounds, not locally probed first dates.
Offsets of 1--7 days represent forecasts made 24--168 hours before valid time,
subject to each model's forecast horizon.\footnote{Open-Meteo,
\href{https://open-meteo.com/en/docs/previous-runs-api#data-availability}{Previous Runs API: Data Availability and API Documentation},
accessed 7 September 2026.}

The local forecast dataset covers \textbf{1--30 June 2026}, with its latest
valid timestamp at \textbf{30 June 2026, 23:00 UTC}, for New York, Los Angeles,
and Miami. Nine of ten requested models returned data; ECMWF AIFS returned
only missing values. Available positive lead offsets reach 7 days for IFS,
GFS, GEM, JMA and NBM; 6 for ICON and UKMO; 3 for M\'et\'eo-France; and 1 for
HRRR. Some combinations have gaps. This local end date is the requested
download boundary, not an API archive cutoff; the newest provider timestamp
has not been measured here.

\textbf{Backtest implication.}
The earliest candidate overlap depends on both city and model: New York/GFS
can potentially extend to August 2021, while a multi-model study generally
starts in 2024 or later. Including Los Angeles moves the market-side start
to 5 January 2025. These are candidate bounds pending forecast and candle
coverage checks. Use positive forecast leads and verify information was
available before the simulated trade; lead zero is diagnostic only.

{\footnotesize
\textbf{Local evidence (paths relative to repository root):}
\path{notebooks_proto/probe_kalshi_daily_high_limits.py};
\path{notebooks_proto/out/kalshi_daily_high_probe/series_summary.json};
\path{data/processed/forecasts/coverage.csv};
\path{data/processed/forecasts/hourly_temp.parquet}.
}
\end{document}
